\chapter{Boolean Analysis — Smolensky Algebra}
%%% generated by the blueprint pipeline from TCSlib.BooleanAnalysis.RazborovSmolensky.SmolenskyAlgebra
%%%%%%%%%%%%%%%%%%%%%%%%%%%%%%%%%%%%%%%%%%%%%%%%%%%%%%%%%%%%%%%%%%%%%%%%%%%%%%%
\section{Overview}

This file develops the Smolensky (algebraic) side of the Razborov--Smolensky lower bound
for $\mathrm{AC}^0[p]$ circuits: it packages low-degree polynomial approximation over
$\mathbb{Z}/p$ on the Boolean cube, converts the pointwise circuit-approximation theorem
into a single low-degree polynomial by averaging, sets up the field
$\mathbb{F}_{p^{q-1}}$ together with a nontrivial $q$-th root of unity $\omega$, and then
works over the root-of-unity cube $\{1,\omega\}^n$: multilinear representation, the split
at degree $n/2$, and the final Hamming-ball counting obstruction.

%%%%%%%%%%%%%%%%%%%%%%%%%%%%%%%%%%%%%%%%%%%%%%%%%%%%%%%%%%%%%%%%%%%%%%%%%%%%%%%
\section{Declarations}

\begin{definition}[MOD $q$ target inside $\mathbb{Z}/p$]
\lean{ACP.modQTarget}
\label{ACP.modQTarget}
\leanok
\uses{ACP.modGateOp}
For a prime $q$, the Boolean function $\mathrm{MOD}_q$ on $n$ inputs viewed as a function
$\{0,1\}^n \to \mathbb{Z}/p$: the output bit of the gate operation $\mathrm{modGateOp}\ q\ n$
is cast into $\mathbb{Z}/p$.
\end{definition}

\begin{definition}[Bad input count on the Boolean cube]
\lean{ACP.badInputCount}
\label{ACP.badInputCount}
\leanok
\uses{ACP.boolInput}
For a target $f : \{0,1\}^n \to \mathbb{Z}/p$ and a polynomial
$P \in (\mathbb{Z}/p)[X_1,\dots,X_n]$, the number of Boolean inputs $x$ with
$P(\mathrm{boolInput}(x)) \neq f(x)$.
\end{definition}

\begin{definition}[Low-degree bad-count lower bound]
\lean{ACP.LowDegreeBadCountLB}
\label{ACP.LowDegreeBadCountLB}
\leanok
\uses{ACP.badInputCount}
The predicate asserting that every polynomial $P$ of total degree at most $d$ has
$E \le \mathrm{badInputCount}(f, P)$, i.e.\ $f$ cannot be computed by a degree-$d$
polynomial on more than $2^n - E$ points of the Boolean cube.
\end{definition}

\begin{lemma}[Averaging a pointwise counting bound over parameters]
\lean{ACP.exists_good_parameter_of_pointwise_bound}
\label{ACP.exists_good_parameter_of_pointwise_bound}
\leanok
\difficulty{5}
Let $\alpha$ and $\beta$ be finite types with $\beta$ nonempty, let $\mathrm{Fail} :
\alpha \to \beta \to \mathrm{Prop}$ be a decidable binary relation, and let $B, C$ be
natural numbers. Suppose that for every $a \in \alpha$ the number of $b \in \beta$ with
$\mathrm{Fail}\,a\,b$ satisfies
\[
\abs{\{\, b \in \beta : \mathrm{Fail}\,a\,b \,\}}\cdot C \;\le\; B\cdot\abs{\beta}.
\]
Then there exists a single $b \in \beta$ for which
\[
\abs{\{\, a \in \alpha : \mathrm{Fail}\,a\,b \,\}}\cdot C \;\le\; B\cdot\abs{\alpha}.
\]
\end{lemma}

\begin{theorem}[A single polynomial from a pointwise error bound]
\lean{ACP.exists_single_polynomial_from_pointwise_distribution}
\label{ACP.exists_single_polynomial_from_pointwise_distribution}
\leanok
\uses{ACP.badInputCount, ACP.boolInput, ACP.exists_good_parameter_of_pointwise_bound}
\difficulty{2}
Fix a prime $p$, and let $(P_s)_{s \in S}$ be a family of polynomials in
$(\bbz/p)[X_1,\dots,X_n]$ indexed by a nonempty finite set $S$ of seeds, let $f :
\{0,1\}^n \to \bbz/p$ be a target function, and let $\ell$ and $B$ be natural numbers.
For a Boolean input $x \in \{0,1\}^n$, write $\bar{x} \in (\bbz/p)^n$ for its
coordinatewise image under the inclusion $\{0,1\} \hookrightarrow \bbz/p$. Suppose that
for every Boolean input $x$,
\[
\bigl|\{\, s \in S : P_s(\bar{x}) \neq f(x)\,\}\bigr| \cdot 2^{\ell} \;\le\; B \cdot
\abs{S}.
\]
Then there exists a seed $s \in S$ for which the number of Boolean inputs $x \in
\{0,1\}^n$ with $P_s(\bar{x}) \neq f(x)$, multiplied by $2^{\ell}$, is at most $B \cdot
2^{n}$.
\end{theorem}

\begin{theorem}[{Low-degree polynomial approximation of an AC⁰[p] circuit}]
\lean{ACP.exists_single_poly_for_circuit_one_size}
\label{ACP.exists_single_poly_for_circuit_one_size}
\leanok
\uses{ACP.ACp_GateOps, ACP.FeedForward, ACP.FeedForward.eval₁, ACP.FeedForward.onlyUsesGates, ACP.FeedForward.size, ACP.badInputCount, ACP.circuitDegreeBound, ACP.exists_poly_distribution_for_circuit_one_size, ACP.exists_single_polynomial_from_pointwise_distribution}
\difficulty{2}
Let $p$ be a prime, let $n\in\bbn$, and let $F$ be a feedforward circuit over the
alphabet $\{0,1\}$ with input type $\mathrm{Fin}\,n$, a unique output node, and finitely
many nodes in each layer, every one of whose gates is drawn from the $\mathrm{AC}^0[p]$
gate set (the identity, NOT, and unbounded fan-in AND gates, together with the unbounded
$\mathrm{MOD}_p$ gates). Write $g\colon\{0,1\}^n\to\bbz/p$ for the function computed by
$F$, namely its single output value cast into $\bbz/p$. Then for every $\ell\in\bbn$
there exists a polynomial $P\in(\bbz/p)[X_1,\dots,X_n]$ of total degree at most
$\bigl((p-1)\ell\bigr)^{\operatorname{depth}(F)}$ such that the number $b$ of Boolean
inputs $x\in\{0,1\}^n$ at which $P$, evaluated at the point of $(\bbz/p)^n$ whose $i$-th
coordinate is the image of $x_i$ in $\bbz/p$, differs from $g(x)$ satisfies
\[
b\cdot 2^{\ell}\;\le\;\operatorname{size}(F)\cdot 2^{n},
\]
where $\operatorname{size}(F)$ is the number of non-input nodes of $F$.
\end{theorem}

\begin{theorem}[Size lower bound from a low-degree bad-count bound]
\lean{ACP.size_lower_bound_from_badCountLB}
\label{ACP.size_lower_bound_from_badCountLB}
\leanok
\uses{ACP.ACp_GateOps, ACP.FeedForward, ACP.FeedForward.eval₁, ACP.FeedForward.onlyUsesGates, ACP.FeedForward.size, ACP.LowDegreeBadCountLB, ACP.badInputCount, ACP.circuitDegreeBound, ACP.exists_single_poly_for_circuit_one_size, ACP.modGateOp, ACP.modQTarget}
\difficulty{3}
Let $p$ and $q$ be primes, and let $F$ be a feedforward circuit over the alphabet
$\{0,1\}$ with input coordinates indexed by $\mathrm{Fin}\,n$, finitely many nodes in
each layer, and a single output node. Suppose every gate of $F$ belongs to the
$\mathrm{AC}^0[p]$ gate set (identity, NOT, unbounded fan-in AND, and unbounded
$\mathrm{MOD}_p$ gates of every arity), and that $F$ computes the $\mathrm{MOD}_q$
function on every Boolean input $x \in \{0,1\}^n$. Let $\ell, E \in \bbn$, and suppose
the low-degree bad-count lower bound holds for $\mathrm{MOD}_q$ (viewed as a function
into $\bbz/p$) with degree budget $\bigl((p-1)\ell\bigr)^{\mathrm{depth}(F)}$ and count
$E$; that is, every polynomial in $(\bbz/p)[X_1,\dots,X_n]$ of total degree at most
$\bigl((p-1)\ell\bigr)^{\mathrm{depth}(F)}$ disagrees with $\mathrm{MOD}_q$ on at least
$E$ points of the Boolean cube. Then
\[
E \cdot 2^{\ell} \le \abs{F} \cdot 2^{n},
\]
where $\abs{F}$ is the number of non-input nodes of $F$.
\end{theorem}

\begin{theorem}[Size lower bound from a relative bad-count bound]
\lean{ACP.size_lower_bound_from_relative_badCountLB}
\label{ACP.size_lower_bound_from_relative_badCountLB}
\leanok
\uses{ACP.ACp_GateOps, ACP.FeedForward, ACP.FeedForward.eval₁, ACP.FeedForward.onlyUsesGates, ACP.FeedForward.size, ACP.LowDegreeBadCountLB, ACP.circuitDegreeBound, ACP.modGateOp, ACP.modQTarget, ACP.size_lower_bound_from_badCountLB}
\difficulty{3}
Let $p$ and $q$ be primes, and let $F$ be a layered feedforward circuit over the Boolean
alphabet with input coordinates $\mathrm{Fin}\,n$ and a single output, each of whose
layers is finite. Suppose $F$ uses only $\mathrm{AC}^0[p]$ gates (the identity gate, the
NOT gate, unbounded fan-in AND gates, and unbounded $\mathrm{MOD}_p$ gates), and that
$F$ computes the $\mathrm{MOD}_q$ function on its $n$ Boolean inputs. Fix $\ell, \delta
\in \bbn$, and assume the relative low-degree bad-count bound: every polynomial $P \in
(\bbz/p)[X_1,\dots,X_n]$ of total degree at most
$\bigl((p-1)\ell\bigr)^{\mathrm{depth}(F)}$ disagrees with the $\mathrm{MOD}_q$ target
(viewed in $\bbz/p$) on at least $\delta \cdot 2^{n}$ of the $2^{n}$ Boolean inputs.
Then the size of $F$ satisfies $\delta \cdot 2^{\ell} \le \abs{F}$.
\end{theorem}

\begin{definition}[Standard field $\mathbb{F}_{p^{q-1}}$]
\lean{ACP.ModqField}
\label{ACP.ModqField}
\leanok
The Galois field $\mathbb{F}_{p^{q-1}}$, the standard field choice for the
$\mathrm{MOD}_q$ lower bound.
\end{definition}

\begin{lemma}[The predecessor of a prime is nonzero]
\lean{ACP.q_sub_one_ne_zero}
\label{ACP.q_sub_one_ne_zero}
\leanok
\difficulty{1}
Let $q$ be a prime number. Then $q - 1 \neq 0$, where the subtraction is taken over the
natural numbers.
\end{lemma}

\begin{lemma}[Cardinality of the field $\mathbb{F}_{p^{q-1}}$]
\lean{ACP.natCard_modqField}
\label{ACP.natCard_modqField}
\leanok
\uses{ACP.ModqField, ACP.q_sub_one_ne_zero}
\difficulty{1}
Let $p$ and $q$ be prime numbers. The Galois field $\mathbb{F}_{p^{q-1}}$ has exactly
$p^{q-1}$ elements.
\end{lemma}

\begin{theorem}[An element of order $q$ in $\mathbb{F}_{p^{q-1}}^{\times}$]
\lean{ACP.exists_unit_of_order_q_modqField}
\label{ACP.exists_unit_of_order_q_modqField}
\leanok
\uses{ACP.ModqField, ACP.natCard_modqField}
\difficulty{5}
Let $p$ and $q$ be distinct prime numbers, and let $\mathbb{F}_{p^{q-1}}$ denote the
finite field with $p^{q-1}$ elements. Then the multiplicative group
$\mathbb{F}_{p^{q-1}}^{\times}$ contains a unit $u$ whose order is exactly $q$.
\end{theorem}

\begin{theorem}[Nontrivial $q$-th root of unity in $\bbf_{p^{q-1}}$]
\lean{ACP.exists_nontrivial_qth_root_modqField}
\label{ACP.exists_nontrivial_qth_root_modqField}
\leanok
\uses{ACP.ModqField, ACP.exists_unit_of_order_q_modqField}
\difficulty{3}
Let $p$ and $q$ be distinct primes. Then the field $\bbf_{p^{q-1}}$ contains an element
$\omega \neq 1$ satisfying $\omega^{q} = 1$; that is, a nontrivial $q$-th root of unity.
\end{theorem}

\begin{definition}[Root-of-unity cube $\{1,\omega\}^n$]
\lean{ACP.rootCube}
\label{ACP.rootCube}
\leanok
The subtype of vectors $x : \mathrm{Fin}\,n \to K$ such that each coordinate satisfies
$x_i = 1$ or $x_i = \omega$.
\end{definition}

\begin{theorem}[Polynomial representative on the root cube]
\lean{ACP.exists_multilinear_representative_on_rootCube}
\label{ACP.exists_multilinear_representative_on_rootCube}
\leanok
\uses{ACP.rootCube}
\difficulty{6}
Let $K$ be a field, let $\omega \in K$, and let $n$ be a natural number. Write
$\{1,\omega\}^n$ for the root cube, the set of vectors $x = (x_1,\dots,x_n) \in K^n$
with $x_i \in \{1,\omega\}$ for every $i$. Then for every function $f : \{1,\omega\}^n
\to K$ there exists a polynomial $P \in K[X_1,\dots,X_n]$ such that $P(x_1,\dots,x_n) =
f(x)$ for all $x$ in the cube.
\end{theorem}

\begin{definition}[Squarefree monomial]
\lean{ACP.squarefreeMonomial}
\label{ACP.squarefreeMonomial}
\leanok
For $s \subseteq \mathrm{Fin}\,n$, the monomial $\prod_{i \in s} X_i$.
\end{definition}

\begin{definition}[Squarefree (multilinear) polynomial from coefficients]
\lean{ACP.squarefreePolynomial}
\label{ACP.squarefreePolynomial}
\leanok
\uses{ACP.squarefreeMonomial}
Given coefficients $c_S$ indexed by subsets $S \subseteq \mathrm{Fin}\,n$, the polynomial
\[
  \sum_{S} c_S \prod_{i \in S} X_i .
\]
\end{definition}

\begin{theorem}[Evaluation of a squarefree monomial]
\lean{ACP.squarefreeMonomial_eval}
\label{ACP.squarefreeMonomial_eval}
\leanok
\uses{ACP.squarefreeMonomial}
\difficulty{1}
Let $K$ be a field, let $n$ be a natural number, and let $s \subseteq \mathrm{Fin}\,n$.
For every point $x : \mathrm{Fin}\,n \to K$, the squarefree monomial $\prod_{i \in s}
X_i$ evaluated at $x$ equals $\prod_{i \in s} x_i$.
\end{theorem}

\begin{theorem}[Affine map realises coordinatewise inversion on the cube]
\lean{ACP.rootCube_affine_inverse}
\label{ACP.rootCube_affine_inverse}
\leanok
\uses{ACP.rootCube}
\difficulty{2}
Let $K$ be a field, let $\omega \in K$ be a nonzero element, and let $x \in
\{1,\omega\}^n$, so that each coordinate $x_i$ equals either $1$ or $\omega$. Then for
every coordinate $i$,
\[
  1 + \omega^{-1} - \omega^{-1} x_i = x_i^{-1}.
\]
\end{theorem}

\begin{theorem}[Total degree of a squarefree monomial]
\lean{ACP.squarefreeMonomial_totalDegree_le_card}
\label{ACP.squarefreeMonomial_totalDegree_le_card}
\leanok
\uses{ACP.squarefreeMonomial}
\difficulty{2}
Fix a field $K$ and let $s \subseteq \{0,1,\dots,n-1\}$ be a set of variable indices.
The total degree of the squarefree monomial $\prod_{i \in s} X_i$ in the polynomial ring
$K[X_0,\dots,X_{n-1}]$ is at most $\abs{s}$.
\end{theorem}

\begin{theorem}[Coordinates of a root-cube vector are nonzero]
\lean{ACP.rootCube_coord_ne_zero}
\label{ACP.rootCube_coord_ne_zero}
\leanok
\uses{ACP.rootCube}
\difficulty{2}
Let $K$ be a field and let $\omega \in K$ with $\omega \neq 0$. If $x \in
\{1,\omega\}^n$, that is, $x$ is a vector in $K^n$ each of whose coordinates equals
either $1$ or $\omega$, then every coordinate is nonzero: $x_i \neq 0$ for each index
$i$.
\end{theorem}

\begin{theorem}[Top monomial times complementary inverse monomial]
\lean{ACP.rootCube_top_mul_compl_inverse}
\label{ACP.rootCube_top_mul_compl_inverse}
\leanok
\uses{ACP.rootCube, ACP.rootCube_coord_ne_zero}
\difficulty{4}
Let $K$ be a field, let $\omega \in K$ be nonzero, and let $x \in \{1,\omega\}^n$, that
is, a vector $x \colon \mathrm{Fin}\,n \to K$ each of whose coordinates equals $1$ or
$\omega$. Then for every subset $s \subseteq \mathrm{Fin}\,n$,
\[
  \Bigl(\prod_{i} x_i\Bigr)\cdot \prod_{i \in s^{c}} x_i^{-1} \;=\; \prod_{i \in s} x_i,
\]
where the first product runs over all coordinates and $s^{c}$ denotes the complement of
$s$.
\end{theorem}

\begin{theorem}[Splitting a multilinear polynomial at half degree on the cube $\{1,\omega\}^n$]
\lean{ACP.split_multilinear_at_half_degree}
\label{ACP.split_multilinear_at_half_degree}
\leanok
\uses{ACP.rootCube, ACP.rootCube_affine_inverse, ACP.rootCube_top_mul_compl_inverse, ACP.squarefreeMonomial, ACP.squarefreeMonomial_totalDegree_le_card, ACP.squarefreePolynomial}
\difficulty{6}
Let $K$ be a field, let $\omega \in K$ be nonzero, and let $(c_S)_{S \subseteq
\mathrm{Fin}\,n}$ be a family of coefficients in $K$ indexed by the subsets of
$\mathrm{Fin}\,n$. Then there exist polynomials $P_1, P_2 \in K[X_i : i \in
\mathrm{Fin}\,n]$, each of total degree at most $\lfloor n/2 \rfloor$, such that for
every point $x \in \{1,\omega\}^n$ (that is, every $x$ with $x_i \in \{1,\omega\}$ for
all $i$),
\[
  \sum_{S \subseteq \mathrm{Fin}\,n} c_S \prod_{i \in S} x_i
   \;=\; P_1(x) + \Bigl(\prod_{i} x_i\Bigr)\, P_2(y),
   \qquad y_i = 1 + \omega^{-1} - \omega^{-1} x_i .
\]
\end{theorem}

\begin{definition}[Affine inverse coordinate polynomial]
\lean{ACP.affineInvPoly}
\label{ACP.affineInvPoly}
\leanok
The degree-one polynomial $1 + \omega^{-1} - \omega^{-1} X_i$, which on $\{1,\omega\}$
with $\omega \neq 0$ computes $x \mapsto x^{-1}$.
\end{definition}

\begin{theorem}[Evaluation of the affine inverse polynomial]
\lean{ACP.affineInvPoly_eval}
\label{ACP.affineInvPoly_eval}
\leanok
\uses{ACP.affineInvPoly}
\difficulty{1}
Let $K$ be a field, let $\omega \in K$, let $n$ be a natural number, and let $i \in
\{0,\dots,n-1\}$ index one of the $n$ coordinates. Then for every point $x =
(x_0,\dots,x_{n-1}) \in K^n$, evaluating the affine inverse coordinate polynomial $1 +
\omega^{-1} - \omega^{-1} X_i$ at $x$ yields \[ 1 + \omega^{-1} - \omega^{-1} x_i. \]
\end{theorem}

\begin{theorem}[Total degree of the affine inverse coordinate polynomial]
\lean{ACP.affineInvPoly_totalDegree_le_one}
\label{ACP.affineInvPoly_totalDegree_le_one}
\leanok
\uses{ACP.affineInvPoly}
\difficulty{3}
Let $K$ be a field, fix an element $\omega \in K$, and let $n$ be a natural number with
a distinguished coordinate index $i \in \{0,\dots,n-1\}$. The polynomial $1 +
\omega^{-1} - \omega^{-1} X_i$ in the polynomial ring $K[X_0,\dots,X_{n-1}]$ has total
degree at most $1$.
\end{theorem}

\begin{definition}[Affine-substituted squarefree monomial]
\lean{ACP.affineSquarefreeMonomial}
\label{ACP.affineSquarefreeMonomial}
\leanok
\uses{ACP.affineInvPoly}
For $s \subseteq \mathrm{Fin}\,n$, the product $\prod_{i \in s}(1 + \omega^{-1} -
\omega^{-1}X_i)$, i.e.\ the squarefree monomial after the affine inverse substitution.
\end{definition}

\begin{theorem}[Evaluation of the affine-substituted squarefree monomial]
\lean{ACP.affineSquarefreeMonomial_eval}
\label{ACP.affineSquarefreeMonomial_eval}
\leanok
\uses{ACP.affineSquarefreeMonomial}
\difficulty{1}
Let $K$ be a field, let $\omega \in K$, and fix $n \in \bbn$. For a subset $s \subseteq
\{1,\dots,n\}$, form the polynomial $\prod_{i \in s}\bigl(1 + \omega^{-1} - \omega^{-1}
X_i\bigr) \in K[X_1,\dots,X_n]$. Then for every point $x \colon \{1,\dots,n\} \to K$,
this polynomial evaluates at $x$ to
\[
\prod_{i \in s}\bigl(1 + \omega^{-1} - \omega^{-1}\, x_i\bigr).
\]
\end{theorem}

\begin{theorem}[Total degree of the affine-substituted squarefree monomial]
\lean{ACP.affineSquarefreeMonomial_totalDegree_le_card}
\label{ACP.affineSquarefreeMonomial_totalDegree_le_card}
\leanok
\uses{ACP.affineInvPoly, ACP.affineInvPoly_totalDegree_le_one, ACP.affineSquarefreeMonomial}
\difficulty{2}
Let $K$ be a field, let $\omega \in K$, and fix $n \in \bbn$ together with a subset $s
\subseteq \{0, 1, \dots, n-1\}$. Then the total degree of the affine-substituted
squarefree monomial $\prod_{i \in s}\bigl(1 + \omega^{-1} - \omega^{-1} X_i\bigr)$,
viewed as a polynomial in $K[X_0, \dots, X_{n-1}]$, is at most $\abs{s}$.
\end{theorem}

\begin{theorem}[Direct half-degree split on the cube $\{1,\omega\}^n$]
\lean{ACP.split_multilinear_at_half_degree_direct}
\label{ACP.split_multilinear_at_half_degree_direct}
\leanok
\uses{ACP.affineSquarefreeMonomial, ACP.affineSquarefreeMonomial_totalDegree_le_card, ACP.rootCube, ACP.rootCube_affine_inverse, ACP.rootCube_top_mul_compl_inverse, ACP.squarefreeMonomial, ACP.squarefreeMonomial_totalDegree_le_card, ACP.squarefreePolynomial}
\difficulty{6}
Let $K$ be a field, fix a scalar $\omega \in K$ with $\omega \neq 0$, and let $n$ be a
natural number. For any family of coefficients $(c_S)$ indexed by the subsets $S
\subseteq \{1,\dots,n\}$, form the multilinear polynomial $F = \sum_{S} c_S \prod_{i \in
S} X_i$ in $K[X_1,\dots,X_n]$. Then there exist polynomials $P_1, R \in
K[X_1,\dots,X_n]$, each of total degree at most $\lfloor n/2 \rfloor$, such that for
every point $x \in \{1,\omega\}^n$ (that is, every $x$ with each coordinate $x_i$ equal
to $1$ or to $\omega$) one has
\[
  F(x) = P_1(x) + \Bigl(\prod_{i=1}^{n} x_i\Bigr)\, R(x).
\]
\end{theorem}

\begin{definition}[Bad count of a polynomial on the root cube]
\lean{ACP.rootCubeBadCount}
\label{ACP.rootCubeBadCount}
\leanok
\uses{ACP.rootCube}
The number of points $x \in \{1,\omega\}^n$ on which $P(x) \neq f(x)$.
\end{definition}

\begin{definition}[Hamming distance between functions on the root cube]
\lean{ACP.rootCubeFunctionBadCount}
\label{ACP.rootCubeFunctionBadCount}
\leanok
\uses{ACP.rootCube}
The number of points $x \in \{1,\omega\}^n$ on which two functions
$f, g : \{1,\omega\}^n \to K$ differ.
\end{definition}

\begin{definition}[Hamming ball on the root cube]
\lean{ACP.rootCubeBall}
\label{ACP.rootCubeBall}
\leanok
\uses{ACP.rootCube, ACP.rootCubeFunctionBadCount}
The finite set of functions $f : \{1,\omega\}^n \to K$ that differ from a given center
function on at most $e$ points of the cube.
\end{definition}

\begin{theorem}[Covering by Hamming balls: a counting bound]
\lean{ACP.finite_cover_by_hamming_balls_card_bound}
\label{ACP.finite_cover_by_hamming_balls_card_bound}
\leanok
\difficulty{5}
Let $\alpha$, $\beta$, and $\mathrm{Cand}$ be finite sets, and assign to each candidate
$c \in \mathrm{Cand}$ a function $\mathrm{center}(c) \colon \alpha \to \beta$. For a
radius $e \in \bbn$, write the Hamming ball of radius $e$ about $c$ for the set of
functions $f \colon \alpha \to \beta$ that differ from $\mathrm{center}(c)$ in at most
$e$ coordinates, that is, with $\abs{\{a \in \alpha : \mathrm{center}(c)(a) \neq f(a)\}}
\le e$. Suppose that every such ball contains at most $B$ functions, and that every $f
\colon \alpha \to \beta$ lies in the ball about at least one candidate. Then \[
\abs{\{\text{functions } \alpha \to \beta\}} \le \abs{\mathrm{Cand}} \cdot B. \]
\end{theorem}

\begin{theorem}[Counting obstruction to low-degree approximation on the root cube]
\lean{ACP.rootCube_counting_obstruction}
\label{ACP.rootCube_counting_obstruction}
\leanok
\uses{ACP.finite_cover_by_hamming_balls_card_bound, ACP.rootCube, ACP.rootCubeBadCount, ACP.rootCubeBall, ACP.rootCubeFunctionBadCount}
\difficulty{4}
Let $K$ be a finite field, fix an element $\omega \in K$, and let $n, D, e, B$ be
natural numbers. Write $\{1,\omega\}^n$ for the cube of vectors in $\mathrm{Fin}\,n \to
K$ whose every coordinate equals $1$ or $\omega$. Suppose a finite index set
$\mathrm{Cand}$ is equipped with a map assigning to each candidate $c$ a polynomial $P_c
\in K[X_1,\dots,X_n]$, and that the following hold: (i) every polynomial $Q$ of total
degree at most $D$ agrees, at every point of $\{1,\omega\}^n$, with $P_c$ for some
candidate $c$; (ii) for each candidate $c$, the Hamming ball of radius $e$ centered at
the function $x \mapsto P_c(x)$ contains at most $B$ functions $\{1,\omega\}^n \to K$;
and (iii) $\abs{\{1,\omega\}^n \to K} > \abs{\mathrm{Cand}}\cdot B$. Then it is not the
case that every function $f : \{1,\omega\}^n \to K$ admits a polynomial $Q$ of total
degree at most $D$ differing from $f$ on at most $e$ points of the cube.
\end{theorem}

\begin{theorem}[Approximating the top monomial approximates every function]
\lean{ACP.rootProd_approx_implies_all_functions_approx}
\label{ACP.rootProd_approx_implies_all_functions_approx}
\leanok
\uses{ACP.rootCube, ACP.rootCubeBadCount, ACP.split_multilinear_at_half_degree_direct, ACP.squarefreePolynomial}
\difficulty{4}
Let $K$ be a finite field and fix $\omega \in K$ with $\omega \neq 0$, and let $n$ be a
natural number. Suppose that every function $f\colon \{1,\omega\}^n \to K$ agrees, at
every point of the cube $\{1,\omega\}^n$, with some squarefree polynomial $\sum_{S
\subseteq \{1,\dots,n\}} c_S \prod_{i \in S} X_i$ over $K$. If a polynomial $P$ in $n$
variables over $K$ has total degree at most $d$ and disagrees with the top monomial $x
\mapsto \prod_{i} x_i$ on at most $e$ points of $\{1,\omega\}^n$, then every function
$f\colon \{1,\omega\}^n \to K$ admits a polynomial $Q$ of total degree at most $\lfloor
n/2 \rfloor + d$ that disagrees with $f$ on at most $e$ points of the cube.
\end{theorem}

\begin{theorem}[No low-degree approximation of the top monomial on the root cube]
\lean{ACP.no_low_degree_rootProd_approx}
\label{ACP.no_low_degree_rootProd_approx}
\leanok
\uses{ACP.rootCube, ACP.rootCubeBadCount, ACP.rootProd_approx_implies_all_functions_approx, ACP.squarefreePolynomial}
\difficulty{2}
Let $K$ be a finite field, let $\omega \in K$ be nonzero, and write $\{1,\omega\}^n$ for
the set of vectors in $K^n$ each of whose coordinates equals $1$ or $\omega$. Suppose
that every function $f : \{1,\omega\}^n \to K$ agrees, at every point of the cube, with
some multilinear polynomial $\sum_{S \subseteq \{1,\dots,n\}} c_S \prod_{i \in S} X_i$;
and suppose that some function on the cube cannot be approximated to within $e$ errors
by any polynomial of total degree at most $\lfloor n/2 \rfloor + d$ — that is, for at
least one $f$, every such polynomial disagrees with $f$ on more than $e$ points of the
cube. Then no polynomial $P$ of total degree at most $d$ agrees with the top monomial $x
\mapsto \prod_{i=1}^{n} x_i$ on all but at most $e$ points of $\{1,\omega\}^n$.
\end{theorem}

\begin{theorem}[Top-monomial inapproximability from finite counting data]
\lean{ACP.no_low_degree_rootProd_approx_of_finite_counting}
\label{ACP.no_low_degree_rootProd_approx_of_finite_counting}
\leanok
\uses{ACP.no_low_degree_rootProd_approx, ACP.rootCube, ACP.rootCubeBadCount, ACP.rootCubeBall, ACP.rootCube_counting_obstruction, ACP.squarefreePolynomial}
\difficulty{1}
Let $K$ be a finite field, let $\omega \in K$ be nonzero, and fix natural numbers $n, d,
e, B$. Suppose that every function $f : \{1,\omega\}^n \to K$ agrees, at every point of
the cube $\{1,\omega\}^n$, with some squarefree (multilinear) polynomial $\sum_{S
\subseteq \{1,\dots,n\}} c_S \prod_{i \in S} X_i$. Suppose further that a finite family
of polynomials $(P_c)_{c \in \mathrm{Cand}}$, indexed by a finite type $\mathrm{Cand}$,
is complete for total degree at most $\lfloor n/2 \rfloor + d$, in the sense that every
$Q \in K[X_1,\dots,X_n]$ with $\deg Q \le \lfloor n/2 \rfloor + d$ agrees with some
$P_c$ at every point of the cube; that for each $c$ the Hamming ball of radius $e$
centered at the function $x \mapsto P_c(x)$ contains at most $B$ functions on the cube;
and that the number of functions $\{1,\omega\}^n \to K$ strictly exceeds
$\abs{\mathrm{Cand}} \cdot B$. Then there is no polynomial $P$ with $\deg P \le d$ such
that $P(x) \neq \prod_{i} x_i$ for at most $e$ points $x$ of the cube; that is, no
polynomial of total degree at most $d$ approximates the product $\prod_i x_i$ on
$\{1,\omega\}^n$ with at most $e$ errors.
\end{theorem}
